% =====================================================================
% Poster template, DIN A0 portrait.
% =====================================================================
\documentclass[final,svgnames,english]{beamer}
\usepackage[english]{babel}

\usepackage[orientation=portrait,size=a0,scale=1.4]{beamerposter}
\usetheme{InriaPoster}

% Headline template without French flag or Inria logo
\setbeamertemplate{headline}{
  \leavevmode
  \begin{beamercolorbox}[wd=\linewidth]{upper separation line head}
    \vskip7.4ex
  \end{beamercolorbox}
  \begin{columns}[b]
    \begin{column}{.02\linewidth}
    \end{column}
    \begin{column}{.96\linewidth}
      \raggedright
      \usebeamercolor{title in headline}{\color{fg}\textbf{\LARGE{\inserttitle}}\\[0.92ex]}
      \usebeamercolor{author in headline}{\color{fg}\large{\insertauthor}\\[0.29ex]}
      \usebeamercolor{institute in headline}{\color{fg}\normalsize{\insertinstitute}\\[0.74ex]}
    \end{column}
    \begin{column}{.02\linewidth}
    \end{column}
  \end{columns}
  \begin{beamercolorbox}[wd=\linewidth]{lower separation line head}
    \vskip2ex
  \end{beamercolorbox}
}

\usepackage{amsmath,amsthm, amssymb, latexsym}
\usepackage{mathspec}
\setmainfont{Inria Sans}[Scale=MatchLowercase, Kerning=Uppercase]
\setsansfont{Inria Sans}[Scale=MatchLowercase, Kerning=Uppercase]
\setmonofont{Fira Code}[Scale=0.74, Kerning=Uppercase]
%\setmathfont(Digits,Latin,Greek,Symbols)[Scale=MatchLowercase]{Inria Serif}
\setmathrm[Scale=MatchLowercase]{Inria Sans}
%\usepackage{mathastext}

% \boldmath %math en gras

\usepackage{graphicx}
\usepackage{booktabs}
\usepackage{qrcode}
\usepackage{adjustbox}
\usepackage{xspace}
\usepackage{nicefrac}
\usepackage{multiobjective}
\usepackage{fontawesome5}
\usepackage{listings}
\lstset{basicstyle=\small\ttfamily,columns=fullflexible,keepspaces=true}
\urlstyle{tt}
\makeatletter
\providecommand{\@ptsize}{0}
\makeatother
\usepackage{setspace}

\usepackage{tikz}
\usetikzlibrary{calc, backgrounds, arrows, shapes}
\usepackage{smartdiagram}
\usepackage{lipsum}

\usepackage[backend=bibtex, style=numeric, autocite=superscript, defernumbers=true, natbib=true]{biblatex}
\usepackage{csquotes}
\usepackage{tikz}           % vector diagrams in native document font
\usetikzlibrary{arrows.meta,positioning,fit,calc}

\tikzset{
  tbox/.style={draw=inria-2024-bleu-nuit, line width=1.5pt, rounded corners=6pt, fill=white, align=center, font=\normalsize, inner sep=6pt, minimum height=2.2cm},
  tbox_optical/.style={tbox, draw=inria-2024-bleu-azur, fill=inria-2024-bleu-azur!10},
  tbox_sar/.style={tbox, draw=inria-2024-bleu-nuit, fill=inria-2024-bleu-nuit!8},
  tbox_model/.style={tbox, draw=inria-2024-violet, fill=inria-2024-violet!10},
  tbox_output/.style={tbox, draw=inria-2024-vert-tendre, fill=inria-2024-vert-tendre!14},
  tbox_loss/.style={tbox, draw=inria-rouge, fill=inria-rouge!10},
  tgroup/.style={draw=inria-2024-gris-bleu!60, line width=1.5pt, rounded corners=8pt, fill=inria-2024-gris-bleu!4, inner sep=14pt},
  tarrow/.style={-{Stealth[length=10pt,width=7pt]}, line width=2pt, draw=inria-2024-bleu-nuit},
  tdash/.style={-{Stealth[length=10pt,width=7pt]}, line width=2pt, dashed, draw=inria-2024-framboise},
  tlbl/.style={font=\small\itshape, fill=white, inner sep=2pt}
}

\makeatletter
\RequireBibliographyStyle{numeric}
%\DeclareFieldFormat{labelnumberwidth}{\mkbibbrackets{#1}}
\DeclareFieldFormat{labelnumberwidth}{#1.}
\makeatother

\DeclareFieldFormat{hal}{%
  \mkbibacro{hal}\addcolon\space
  \ifhyperref
  {\href{https://hal.archives-ouvertes.fr/#1}{\nolinkurl{#1}}}
{\nolinkurl{#1}}}

\DeclareFieldAlias{eprint:hal}{hal}
\DeclareFieldAlias{eprint:HAL}{eprint:hal}

\renewbibmacro*{eprint}{%
  \printfield{hal}%
  \newunit\newblock
  \iffieldundef{eprinttype}
  {\printfield{eprint}}
{\printfield[eprint:\strfield{eprinttype}]{eprint}}}

\DeclareFieldFormat{url}{%
  \mkbibacro{doi}\addcolon\space
  \ifhyperref
  {\href{https://doi.org/#1}{\nolinkurl{#1}}}
{\nolinkurl{#1}}}

\DeclareFieldFormat{url}{url\addcolon\space\url{#1}}

\renewbibmacro*{url+urldate}{%
  \ifthenelse{\iffieldundef{url}\OR\NOT\iffieldundef{doi}}
  {}
  {\iffieldundef{urlyear}
    {}
    {%\bibstring{retrieved}%
      \setunit{\addspace}%
      %\printurldate
      %\setunit{\urldatecomma}%
      %\bibstring{from}%
      %\setunit{\addspace}
    }%
\iffieldundef{url}{}{\printfield{url}\renewcommand*{\finentrypunct}{\relax}}}}

% Update self.bib with your own publication information
% DO NOT change the "this-paper" key inside self.bib.
\addbibresource{self.bib} % Bibliography file
\addbibresource{references.bib}

% Avoid that our paper shows up in bibliography
\DeclareBibliographyCategory{thispapercat}
\addtocategory{thispapercat}{this-paper}
\AtBeginBibliography{\onehalfspacing}




% ---------------------------------------------------------------------
% Highlight boxes
% ---------------------------------------------------------------------
\newtcolorbox{callout}{%
  enhanced, sharp corners, colback=inria-2024-rouge!11, width=\linewidth,
  colframe=inria-2024-rouge, boxrule=0pt, leftrule=7.4pt,
  left=11pt, right=5pt, top=9pt, bottom=9pt,
before skip=11pt, after skip=5pt}

\newtcolorbox{takeaway}{%
  enhanced, sharp corners, colback=inria-2024-vert-tendre!11, width=\linewidth,
  colframe=inria-2024-vert-tendre, boxrule=0pt, leftrule=7.4pt,
  left=11pt, right=5pt, top=9pt, bottom=9pt,
before skip=11pt, after skip=5pt}

\newtcolorbox{inform}{%
  enhanced, sharp corners, colback=inria-2024-gris-bleu!11, width=\linewidth,
  colframe=inria-2024-gris-bleu, boxrule=0pt, leftrule=7.4pt,
  left=11pt, right=5pt, top=9pt, bottom=9pt,
before skip=11pt, after skip=5pt}

% block for the main findings of the paper
\newtcolorbox{findingblock}[2]{%
  enhanced, sharp corners, colback=inria-2024-#2!1, width=\linewidth,
  colframe=inria-2024-#2, boxrule=0.47pt, leftrule=7.4pt,title={#1}, detach title,
  before upper={\textbf{\tcbtitle}\ },
  coltitle=inria-2024-#2,
  left=11pt, right=5pt, top=9pt, bottom=9pt,
before skip=5pt, after skip=5pt}

\title{On the Complementarity of S-band and C-band SAR Images for Snow Cover Segmentation}
\author{R. Virinchi 254161009 \url{r.virinchi@iitg.ac.in}}
\institute[]{Indian Institute of Technology Guwahati}

\begin{document}
\begin{frame}[fragile]{}
  \vspace{-2em}
  \begin{columns}[T]
    % Forcing left margin
    \begin{column}{0.01\linewidth}
    \end{column}

    % Content
    \begin{column}{0.98\linewidth}
      % TL;DR section: continuous gradient line spanning the whole page above the sentences
      \noindent\begin{tikzpicture}
        \shade[inria degrade, rounded corners=2pt] (0,0) rectangle (\linewidth, \filetheight);
      \end{tikzpicture}\par
      \vspace{0.4ex}
      \begin{columns}[t]
        \begin{column}{0.32\linewidth}
          \begin{plainblock}{}
            \large Snow cover estimation has always been a difficult problem to solve...\raggedright
          \end{plainblock}
        \end{column}%\hfill
        \begin{column}{0.32\linewidth}
          \begin{plainblock}{}
            \large \ldots First there were optical based methods, \tcblower \large \textbf{but they didn't work when there were clouds!}
          \end{plainblock}
        \end{column}%\hfill
        \begin{column}{0.32\linewidth}
          \begin{plainblock}{}
            \large Hence SAR based methods were created... \textcolor{inria-rouge}{\textbf{This work proposes one such method!}}
          \end{plainblock}
        \end{column}
      \end{columns}
      % Main content
      \begin{columns}[T]
        \begin{column}{.485\linewidth}
          \begin{inriastartblock}{Context}
            \begin{itemize}
              \item Nothing can beat snow cover estimated using optical images, when the sky is clear.
              \item But, sky is not always clear. During winter $>50\%$ of sky over a region can be filled with clouds.
              \item Synthetic Aperture Radar is an active remote sensing technique that can penetrate clouds!
              \item But it cannot identify snow on its own.
              \item Enter Deep Learning - We can teach it to identify snow!
            \end{itemize}
        \end{inriastartblock}
        \begin{inriablock}{Weakly Supervised Framework}
            \begin{figure}[!htbp]
              \centering
              \begin{tikzpicture}[node distance=0.85cm and 1.2cm]
                % -------------------------------------------------------------
                % 1. Training Phase (Clear-Sky Days)
                % -------------------------------------------------------------
                % Row 1: Optical weak label branch
                \node[tbox_optical, text width=6.6cm, minimum height=1.9cm] (s2) {\textbf{Sentinel-2}\\[2pt]\small(Optical Image)};
                \node[tbox_optical, text width=6.6cm, minimum height=1.9cm, right=1.2cm of s2] (lis) {\textbf{Snow Labels}\\[2pt]\small(Let-It-Snow)};
                \draw[tarrow] (s2) -- (lis);

                % Row 2: SAR forward pass branch
                \node[tbox_sar, text width=6.6cm, minimum height=1.9cm, below=0.85cm of s2] (s1) {\textbf{Sentinel-1}\\[2pt]\small(SAR Backscatter)};
                \node[tbox_model, text width=6.6cm, minimum height=1.9cm, right=1.2cm of s1] (unet) {\textbf{U-Net Model}\\[2pt]\small(Segmentation)};
                \node[tbox_output, text width=6.6cm, minimum height=1.9cm, right=1.2cm of unet] (pred) {\textbf{Predicted}\\[2pt]\textbf{Snow Map}};
                \node[tbox_loss, text width=6.6cm, minimum height=1.9cm, right=1.2cm of pred] (loss) {\textbf{Loss Function}\\[2pt]\small(Cross-Entropy)};

                \draw[tarrow] (s1) -- node[above, font=\small\itshape] {Input} (unet);
                \draw[tarrow] (unet) -- (pred);
                \draw[tarrow] (pred) -- (loss);

                % Connect LIS weak labels to Loss Calculation (supervision target)
                \draw[tdash] (lis.east) -| node[pos=0.35, above, font=\small\bfseries, text=inria-2024-framboise] {Weak Labels} (loss.north);

                % Training loop feedback: backpropagation to U-Net
                \coordinate (loopb) at ($(loss.south)+(0,-0.65cm)$);
                \draw[tarrow, draw=inria-rouge] (loss.south) -- (loopb) -- (loopb -| unet.south) node[midway, below=2pt, font=\small\bfseries, text=inria-rouge] {Backpropagation (Training Loop)} -- (unet.south);

                % Group box for Training Phase (fully enclosing feedback loop and label)
                \node[tgroup, fit=(s2)(s1)(lis)(unet)(pred)(loss)(loopb), inner xsep=12pt, inner ysep=8pt, label={[font=\normalsize\bfseries, text=inria-2024-bleu-nuit]above:\textbf{1. Training Phase (Clear-Sky Days)}}] (gtrain) {};

                % -------------------------------------------------------------
                % 2. Inference Phase (All-Weather / Cloudy)
                % -------------------------------------------------------------
                \node[tbox_model, text width=7.2cm, minimum height=1.9cm, below=1.2cm of gtrain.south] (unett) {\textbf{Trained U-Net}\\[2pt]\small(Frozen Weights)};
                \node[tbox_sar, text width=7.2cm, minimum height=1.9cm, left=1.4cm of unett] (s1t) {\textbf{Sentinel-1}\\[2pt]\small(Cloudy / All-Weather)};
                \node[tbox_output, text width=7.2cm, minimum height=1.9cm, right=1.4cm of unett] (out) {\textbf{Snow Cover Map}\\[2pt]\small(Segmented Output)};

                \draw[tarrow] (s1t) -- node[above, font=\small\itshape] {Input} (unett);
                \draw[tarrow] (unett) -- (out);

                \node[tgroup, fit=(s1t)(unett)(out), inner xsep=12pt, inner ysep=8pt, label={[font=\normalsize\bfseries, text=inria-2024-bleu-nuit]above:\textbf{2. Inference Phase (All-Weather / Cloudy)}}] (ginf) {};
              \end{tikzpicture}
              \vspace{0.2em}
              \caption{Weakly Supervised Training and Inference Flow Diagram}
              \label{fig:concept}
            \end{figure}
          \end{inriablock}

          \begin{inriablock}{Defining the Objectives:}
            \begin{enumerate}
              \item Build reusable satellite image preprocessing pipeline. ($70\%$ of the work 'cause satellite images contain too much information!)
              \item Generate labels needed to train the model using optical images (hence weakly supervised).
              \item Train model on C-band and get elevation wise metrics.
              \item Train model on S-band and get elevation wise metrics.
              \item Prove they are complementary and device a novel approach to combine them!
              \item Use the new model to outperform previous best-performing model of \citet{briand2026weakly}.
              \item Package everything into a GUI for demonstration.
            \end{enumerate}
            % \begin{takeaway}
              
            % \end{takeaway}
          \end{inriablock}

          \begin{inriablock}{Laying out the Timeline}
            \begin{itemize}
                \item Result by end of February.
                  \begin{itemize}
                    \item Either positive or negative result. Both are meaningful!
                  \end{itemize}
                \item Submit paper by end of March.
                \item Create GUI by end of April.
                  \begin{itemize}
                    \item Time taken $\propto$ no.of external dependencies.
                    \item Because they have a tendency to not cooperate.
                  \end{itemize}
              \end{itemize}
          \end{inriablock}

          % % a simple block to put context
          % \begin{plainblock}{}
          %   \centering
          %   %\begin{adjustbox}{max width=0.83\textwidth}
          %   % Adapted from https://texample.net/3d-cone/
\begin{tikzpicture}[join=round,scale=3.8]
    \tikzstyle{conefill} = [fill=inria-2024-bleu-canard!20,fill opacity=0.8]
    \tikzstyle{ann} = [fill=white,font=\footnotesize,inner sep=1pt]
    \tikzstyle{ghostfill} = [fill=white]
         \tikzstyle{ghostdraw} = [draw=black!50]
    \filldraw[conefill](-.775,1.922)--(-1.162,.283)--(-.274,.5)
                        --(-.183,2.067)--cycle;
    \filldraw[conefill](-.183,2.067)--(-.274,.5)--(.775,.424)
                        --(.516,2.016)--cycle;
    \filldraw[conefill](.516,2.016)--(.775,.424)--(1.369,.1)
                        --(.913,1.8)--cycle;
    \filldraw[conefill](-.913,1.667)--(-1.369,-.1)--(-1.162,.283)
                        --(-.775,1.922)--cycle;
    \draw(1.461,.107)--(1.734,.127);
    \draw[arrows=<->](1.643,1.853)--(1.643,.12);
    \filldraw[conefill](.913,1.8)--(1.369,.1)--(1.162,-.283)
                        --(.775,1.545)--cycle;
    \draw[arrows=->,line width=.4pt](.274,-.5)--(0,0)--(0,2.86);
    \draw[arrows=-,line width=.4pt](0,0)--(-1.369,-.1);
    \draw[arrows=->,line width=.4pt](-1.369,-.1)--(-2.1,-.153);
    \filldraw[conefill](-.516,1.45)--(-.775,-.424)--(-1.369,-.1)
                        --(-.913,1.667)--cycle;
    \draw(-1.369,.073)--(-1.369,2.76);
    \draw(1.004,1.807)--(1.734,1.86);
    \filldraw[conefill](.775,1.545)--(1.162,-.283)--(.274,-.5)
                        --(.183,1.4)--cycle;
    \draw[arrows=<->](0,2.34)--(-.913,2.273);
    \draw(-.913,1.84)--(-.913,2.447);
    \draw[arrows=<->](0,2.687)--(-1.369,2.587);
    \filldraw[conefill](.183,1.4)--(.274,-.5)--(-.775,-.424)
                        --(-.516,1.45)--cycle;
    \draw[arrows=<-,line width=.4pt](.42,-.767)--(.274,-.5);
    \node[ann] at (-.456,2.307) {$r_0$};
    \node[ann] at (-.685,2.637) {$r_1$};
    \node[ann] at (1.643,.987) {$h$};
    \path (.42,-.767) node[below] {$x$}
        (0,2.86) node[above] {$y$}
        (-2.1,-.153) node[left] {$z$};
    % Second version of the cone
    \begin{scope}[xshift=3.5cm]
    \filldraw[ghostdraw,ghostfill](-.775,1.922)--(-1.162,.283)--(-.274,.5)
                                   --(-.183,2.067)--cycle;
    \filldraw[ghostdraw,ghostfill](-.183,2.067)--(-.274,.5)--(.775,.424) 
                                   --(.516,2.016)--cycle;
    \filldraw[ghostdraw,ghostfill](.516,2.016)--(.775,.424)--(1.369,.1)
                                   --(.913,1.8)--cycle;
    \filldraw[ghostdraw,ghostfill](-.913,1.667)--(-1.369,-.1)--(-1.162,.283)
                                   --(-.775,1.922)--cycle;
    \filldraw[ghostdraw,ghostfill](.913,1.8)--(1.369,.1)--(1.162,-.283)
                                   --(.775,1.545)--cycle;
    \filldraw[ghostdraw,ghostfill](-.516,1.45)--(-.775,-.424)--(-1.369,-.1)
                                   --(-.913,1.667)--cycle;
    \filldraw[ghostdraw,ghostfill](.775,1.545)--(1.162,-.283)--(.274,-.5)
                                   --(.183,1.4)--cycle;
    \filldraw[fill=inria-rouge,fill opacity=0.5](-.516,1.45)--(-.775,-.424)--(.274,-.5)
                                         --(.183,1.4)--cycle;
    \fill(-.775,-.424) circle (2pt);
    \fill(.274,-.5) circle (2pt);
    \fill(-.516,1.45) circle (2pt);
    \fill(.183,1.4) circle (2pt);
    \path[font=\footnotesize]
            (.913,1.8) node[right] {$i\hbox{$=$}0$}
            (1.369,.1) node[right] {$i\hbox{$=$}1$};
    \path[font=\footnotesize]
            (-.645,.513) node[left] {$j$}
            (.228,.45) node[right] {$j\hbox{$+$}1$};
    \draw (-.209,.482)+(-60:.25) [yscale=1.3,->] arc(-60:240:.25);
    \fill[black,font=\footnotesize]
                    (-.516,1.45) node [above] {$P_{00}$}
                    (-.775,-.424) node [below] {$P_{10}$}
                    (.183,1.4) node [above] {$P_{01}$}
                    (.274,-.5) node [below] {$P_{11}$};
    \end{scope}
\end{tikzpicture}
          %   %\end{adjustbox}
          % \end{plainblock}

          \begin{inriablock}{Poster Design Principles}
            \begin{itemize}
              \item \textbf{Minimalist:} less is more; focus on the readability of the content.
              \item \textbf{Sophisticated:} clean styling and visual balance.
              \item \textbf{Structured:} organized blocks with clear visual hierarchy.
                \begin{itemize}
                  \item Highlight key takeaways and quantitative metrics.
                  \item Maintain high contrast and legible typography.
                \end{itemize}
            \end{itemize}
          \end{inriablock}

          \begin{inriablock}{The (infinite) NLP loop}
            %\vspace{-2ex}
            \begin{columns}[T]
              \begin{column}{0.5\textwidth}
                \begin{itemize}
                  \item Be concise.
                    \begin{itemize}
                      \item Avoid long narrative text spans,
                      \item instead, itemize or enumerate.
                    \end{itemize}
                  \item Be precise.
                  \item Be illustrative.
                  \item Be self-contained.
                  \item Ensure readability at a distance.
                  \item Stick to a consistent color palette.
                    \begin{itemize}
                      \item you earn extra points if your plots also use those colors!
                    \end{itemize}
                \end{itemize}
              \end{column}
              \begin{column}{0.45\textwidth}
                \centering
                \tikzset{every shadow/.style={fill=none,shadow scale=0}}
\tikzset{module/.append style={shade=none, fill=\col}}
\smartdiagramset{
    border color=\col,
    text color=white,
    module shape=rounded rectangle,
    font=\small,
    text width=0.2\linewidth,
    set color list={inria-2024-framboise,inria-2024-violet,inria-2024-bleu-nuit,inria-2024-bleu-canard,inria-2024-bleu-azur,inria-2024-gris-bleu, inria-2024-cactus,inria-2024-orange},
    circular final arrow disabled=true,
    circular distance=0.38\linewidth,
    arrow tip=to,
    arrow line width=4pt}
\smartdiagram[circular diagram:clockwise]{Paper \texttt{.tex}, Venue rules, Prepare poster, Print, Pin to wall, Present, Take photo, Amaze friends!}
              \end{column}
            \end{columns}
          \end{inriablock}

          % You can move this block around to fit your poster content
          \begin{plainblock}{}
            \begin{columns}[c]
              \begin{column}{0.092\textwidth}
                \qrcode[height=\textwidth,hyperlink]{https://github.com/tensor-calculus/Snow-Cover-Estimation}
              \end{column}
              \begin{column}{0.9\textwidth}%
                {\linespread{0.83}\footnotesize\fullcite{this-paper}\\[-1.28ex]
                \textbf{Code:} \url{https://github.com/tensor-calculus/Snow-Cover-Estimation}}
              \end{column}
            \end{columns}
          \end{plainblock}
        \end{column}

        % Column 2
        \begin{column}{.485\linewidth}
          \begin{plainblock}{}
            \begin{findingblock}{Finding 1. Using \texttt{findingblock} to draw attention.}{framboise}

              You can use \texttt{findingblock} to draw attention to the main findings of the paper.
                \begin{lstlisting}
\begin{findingblock}{<Title of finding>}{<color>}
    <text of the finding>
\end{findingblock}
                \end{lstlisting}
              \begin{itemize}
                \item \textbf{\texttt{color}}: rouge, framboise, violet, bleu-nuit, bleu-canard, bleu-azur, bleu-vert, gris-bleu, cactus, vert-tendre, jaune, orange, and sable.
                \item \textbf{\texttt{Title of finding}}: title of the finding, can be empty.
                \item \textbf{Note:} This is a nice addition to the poster template.
              \end{itemize}
            \end{findingblock}

            \begin{findingblock}{Finding 2. This is actually cool.}{violet}
              \begin{columns}[T]
                \begin{column}{0.6\textwidth}
                  \begin{itemize}
                    \item \lipsum[1][1]
                    \item \lipsum[1][2]
                    \item \lipsum[1][3]
                    \item \lipsum[1][4]
                  \end{itemize}
                \end{column}
                \begin{column}{0.38\textwidth}
                      \begin{tikzpicture}[scale=2.9, every node/.style={font=\small}]
    % Define coordinates for the centers of the 3 circles
    \coordinate (A) at (0, 0.6);
    \coordinate (B) at (-0.7, -0.4);
    \coordinate (C) at (0.7, -0.4);

    % Draw and fill the circles with opacity so intersections show up
    \draw[thick, fill=inria-2024-rouge, fill opacity=0.74] (A) circle (1.2cm);
    \draw[thick, fill=inria-2024-gris-bleu, fill opacity=0.74] (B) circle (1.2cm);
    \draw[thick, fill=inria-2024-framboise, fill opacity=0.74] (C) circle (1.2cm);

    % Outer Set Labels
    \node[color=white] at (0, 1.4) {Research};
    \node[color=white] at (-1.2, -0.92) {\LaTeX};
    \node[color=white] at (1.2, -0.92) {Poster};

    % 1. Label for the exact center overlap (A intersect B intersect C)
    % Uses barycentric coordinates to find the exact middle of A, B, and C
    \node[text opacity=1, color=white] at (barycentric cs:A=1,B=1,C=1) {You};

    % 2. Labels for dual-overlaps (e.g., just A and B)
    %\node[text opacity=1] at (barycentric cs:A=1,B=1) {$A \cap B$};
    %\node[text opacity=1] at (barycentric cs:A=1,C=1) {$A \cap C$};
    %\node[text opacity=1] at (barycentric cs:B=1,C=1) {$B \cap C$};
\end{tikzpicture}
                \end{column}
              \end{columns}
            \end{findingblock}

            \begin{findingblock}{Finding 3. We have other blocks.}{framboise}
              \begin{takeaway}
                \texttt{takeaway}: green block for good news. \lipsum[1][1-2]
              \end{takeaway}

              \begin{callout}
                \texttt{callout}: a red block to draw attention. \lipsum[1][1-2]
              \end{callout}

              \begin{inform}
                \texttt{inform}: to frame some important information or definition. \lipsum[1][1-2]
              \end{inform}
            \end{findingblock}

            \begin{findingblock}{Finding 4. Columns inside your findings are your friend.}{violet}
              \begin{columns}[T]
                \begin{column}{0.47\textwidth}
                  \begin{itemize}
                    \item You can organize your content as (nested) columns.
                      \begin{itemize}
                        \item helps to have narrative next to visual elements,
                        \item it can be fun.
                      \end{itemize}
                  \end{itemize}
                \end{column}
                \begin{column}{0.47\textwidth}
                  \centering
                  % taken from https://texample.net/polygon-division/

\def\hepta{\draw(A) -- (B) -- (C) -- (D) -- (E) -- (F) -- (G) -- cycle;}

% Macro for drawing polygon diagonals. 
% Example \slice{A/C,C/E,E/G,C/G}
\newcommand{\slice}[1]{%
    \hepta
    \draw \foreach \x/\y in {#1} {(\x)--(\y)};
}


\begin{tikzpicture}
    % Define the heptagon coordinates
    \coordinate (A) at (-0.76,1.54);
    \coordinate (B) at (-0.76,0.69);
    \coordinate (C) at (-0.10,0.16);
    \coordinate (D) at (0.73,0.35);
    \coordinate (E) at (1.1,1.11);
    \coordinate (F) at (0.73,1.88);
    \coordinate (G) at (-0.10,2.07);

\matrix[column sep=0.8cm,row sep=0.5cm]
{
    \slice{A/C,C/E,E/G,C/G}&
    \slice{A/C,C/E,E/G,A/E}&
    \slice{A/C,C/E,A/E,A/F}&
    \slice{A/C,C/E,C/F,A/F}&
    \slice{A/C,C/E,C/F,C/G}&
    \slice{A/C,C/F,D/F,A/F}\\
    \slice{A/C,C/F,D/F,C/G}&
    \slice{A/C,C/G,D/G,D/F}&
    \slice{A/C,C/G,D/G,E/G}&
    \slice{A/C,A/D,D/F,A/F}&
    \slice{A/C,A/D,D/F,D/G}&
    \slice{A/C,A/D,D/G,E/G}\\
    \slice{A/C,A/D,A/E,E/G}&
    \slice{A/C,A/D,A/E,A/F}&
    \slice{A/D,B/D,D/F,A/F}&
    \slice{A/D,B/D,D/F,D/G}&
    \slice{A/D,B/D,D/G,E/G}&
    \slice{A/D,B/D,A/E,E/G}\\
    \slice{A/D,B/D,A/E,A/F}&
    \slice{A/E,B/E,B/D,E/G}&
    \slice{A/E,B/E,B/D,A/F}&
    \slice{A/E,B/E,C/E,E/G}&
    \slice{A/E,B/E,C/E,A/F}&
    \slice{A/F,B/F,B/D,D/F}\\
    \slice{A/F,B/F,B/D,B/E}&
    \slice{A/F,B/F,B/E,C/E}&
    \slice{A/F,B/F,C/F,C/E}&
    \slice{A/F,B/F,C/F,D/F}&
    \slice{B/G,B/D,D/F,B/F}&
    \slice{B/G,B/D,D/F,D/G}\\
    \slice{B/G,B/D,D/G,E/G}&
    \slice{B/G,B/D,B/E,E/G}&
    \slice{B/G,B/D,B/E,B/F}&
    \slice{B/G,B/E,C/E,E/G}&
    \slice{B/G,B/E,C/E,B/F}&
    \slice{B/G,B/F,C/F,C/E}\\
    \slice{B/G,B/F,C/F,D/F}&
    \slice{B/G,C/G,C/E,E/G}&
    \slice{B/G,C/G,C/E,C/F}&
    \slice{B/G,C/G,C/F,D/F}&
    \slice{B/G,C/G,D/G,D/F}&
    \slice{B/G,C/G,D/G,E/G}\\
};
\end{tikzpicture}
                \end{column}
              \end{columns}
            \end{findingblock}
          \end{plainblock}

          \begin{inriafinalblock}{Final remarks}
            \begin{itemize}
              \item Use \verb|\begin{inriafinalblock}{header}...\end{inriafinalblock}| to end your poster.
              \item \textbf{Project repository:}
                \begin{itemize}
                  \item \url{https://github.com/tensor-calculus/Snow-Cover-Estimation}
                \end{itemize}
              \item \textbf{Structure \& Clarity:}
                \begin{itemize}
                  \item Keep text concise and prefer visual figures and diagrams.
                  \item High-contrast blocks ensure readability from a distance.
                \end{itemize}
            \end{itemize}
          \end{inriafinalblock}
          % In general, it does not make sense to include a bibliography, but here you have it in case you need it.
          \begin{plainblock}{\small References}
            %%\renewcommand*{\bibfont}{\footnotesize}
            \setlength{\bibitemsep}{0.29em}
            \printbibliography[notcategory=thispapercat]
          \end{plainblock}
        \end{column}
      \end{columns}
    \end{column}

    % right margin
    \begin{column}{0.01\linewidth}
    \end{column}
  \end{columns}
\end{frame}
\end{document}
